\chapter{Intelligence artificielle et assistance à l'analyse}

\section{Objectif de l'assistant IA}
L'assistant IA n'est pas le produit final de SRS~: c'est un intervieweur intelligent
dont le rôle est de collecter des informations. Le produit réel est la transformation
des réponses de l'entretien en modèle de projet structuré, puis en exigences, en
architecture et en documents.

\section{Fonctionnement de l'entretien intelligent}
À chaque tour, l'IA reçoit un contexte structuré (concept courant, historique des
échanges, informations déjà connues, informations manquantes) plutôt qu'une simple
conversation non structurée, afin de générer une question courte, précise et non
redondante.

\section{Gestion du contexte}
Le contexte transmis à l'IA distingue explicitement les informations connues du projet,
les hypothèses, et les questions nécessitant une confirmation — évitant que l'IA
présente une supposition comme un fait acquis.

\section{Restrictions du domaine de l'IA}
L'IA est strictement cantonnée aux sujets liés au cadrage de la plateforme (besoins,
architecture, données, sécurité, déploiement, tests). Toute question hors périmètre
reçoit un message de recentrage explicite plutôt qu'une réponse de chatbot généraliste.

\section{Gestion des fournisseurs IA}
Le client IA est conçu derrière une interface (\texttt{AIClient}) découplée du reste de
l'application, afin de pouvoir ajouter ou remplacer un fournisseur sans modifier le
moteur d'entretien lui-même.

\section{Fournisseurs IA intégrés}
\TODO{Lister ici, avec exactitude, l'état réel de chaque fournisseur au moment de la
rédaction finale (implémenté et testé / partiellement implémenté / non implémenté),
plutôt que l'ambition initiale du cahier des charges.}

\section{Sécurité des clés API}
Les clés API des fournisseurs IA ne transitent jamais côté client~: elles sont lues
côté serveur à partir de variables d'environnement, ne sont jamais journalisées, ni
incluses dans les documents générés, ni committées dans le dépôt Git.

\section{Limites et risques liés à l'IA}
\begin{itemize}
    \item Une réponse d'IA reste une interprétation~: elle doit toujours être
    distinguée d'une information officiellement confirmée par une partie prenante.
    \item En cas d'indisponibilité du fournisseur IA, le Classic Questionnaire reste
    pleinement fonctionnel afin de ne jamais bloquer un entretien.
    \item Aucune information confidentielle non confirmée par la Fondation OCP ne doit
    être présentée comme un fait acquis dans les documents générés.
\end{itemize}
